\chapter{Continuous batching and precision}

Paged KV solves a storage problem; it does not decide who runs next. A serving engine receives prompts of different lengths at different times while existing requests expect regular next-token progress. A second pressure appears simultaneously: KV and model traffic consume memory bandwidth, encouraging lower precision. Both policies can improve resource efficiency and both can damage user-visible quality if treated as slogans.

This chapter makes the policies explicit.

\section{Iteration-level scheduling}
Autoregressive generation is multi-iteration. A request that produces 200 output tokens re-enters the model roughly 200 decode iterations after prefill. Request-level static batching waits too long to change membership.

Orca's iteration-level scheduling formalized the idea that finished requests can leave and new work can enter between iterations. This chapter adopts the concept but defines its own teaching policy.

\section{Resource-aware admission}
Before scheduling, admission verifies model availability and enough KV/page capacity under Chapter 36's allocator policy.
\Plate{01-admission}{A scheduler cannot repair an admission promise that exceeds physical state capacity.}{evod37:admission}

Admission and scheduling remain different: admission says work may enter; scheduling chooses the next iteration's work.

\section{Prefill and decode are different}
A prompt prefill can contain hundreds or thousands of tokens. A decode request typically needs one new token step. Combining an arbitrarily long prefill with ongoing decodes can create large iteration-time variance and poor token latency.

Treating both as simply ``one request'' hides the work asymmetry.

\section{Chunked prefill}
Split a prompt into chunks of at most $C$ tokens. With iteration token budget $B$, schedule
\[
n_{\rm prefill}\le \min(C,B_{\rm remaining}).
\]
\Plate{02-mix}{The teaching scheduler services ready decodes then fills remaining budget with bounded prefill chunks.}{evod37:mix}

Decode-first is an explicit policy choice here, not a theorem. Prefill-prioritizing or balanced policies may produce better throughput or fairness under other SLOs.

\section{Continuous batch churn}
After each iteration:
\begin{enumerate}
\item mark completed decodes;
\item release/request cache pages as appropriate;
\item admit eligible waiting requests;
\item form the next mixed iteration.
\end{enumerate}
A request's causal order is never changed even as batch membership changes.

\section{TTFT}
For arrival time $a_r$ and first emitted output time $f_r$:
\[
\operatorname{TTFT}_r=f_r-a_r.
\]
Time to first token includes queueing, prefill and the first decode. A scheduler that protects ongoing decodes can increase a new long prompt's TTFT.

\section{Inter-token latency}
For output timestamps $t_1,\ldots,t_n$:
\[
\operatorname{ITL}_i=t_{i+1}-t_i.
\]
\Plate{03-metrics}{TTFT and ITL expose different phases of service.}{evod37:metrics}

Report mean only with distributions/tails. Interactive quality is often dominated by p95/p99 stalls, not average token spacing.

\section{Throughput is another axis}
Throughput can be measured as useful generated tokens/s, total processed tokens/s or requests/s. Name the numerator.

A scheduler that increases total prompt-token throughput while causing long decode stalls may be unacceptable for interactive SLOs.

\Plate{06-frontier}{Serving systems operate on a throughput--latency frontier rather than at one universally optimal point.}{evod37:frontier}

\section{Fairness}
Work conservation does not imply fairness. Large prompts, long generations or clients with many concurrent requests can dominate service under naive policies.

Recent serving research explicitly studies fairness under continuous batching. A production scheduler should state its unit of fairness: request, token work, client, weighted entitlement or deadline/SLO class.

\section{Precision is typed}
``INT8 model'' is incomplete. Distinguish:
\[
P=(P_W,P_A,P_{KV},P_{\rm accum},P_{\rm out}).
\]
\Plate{04-precision}{Weights, activations, KV and accumulators have separate precision contracts.}{evod37:precision}

A system may use low-precision storage while accumulating a reduction at higher precision.

\section{FP16 and BF16}
Both occupy 16 bits but trade exponent and mantissa differently. BF16's wider exponent range can tolerate large magnitudes that overflow FP16 more easily; FP16 has greater significand precision over values it can represent.

The right choice depends on model/kernel/hardware and numerical behavior. This chapter's Python conversions study representation error only.

\section{FP8}
FP8 families use eight-bit floating formats with different exponent/mantissa allocations and often scaling metadata. Hardware such as modern accelerators can provide fast FP8 matrix operations, but speed depends on supported kernels and shape/layout.

Do not infer an FP8 speedup from the bit width alone.

\section{INT8}
For symmetric per-tensor INT8 quantization:
\[
s=\frac{\max_i|x_i|}{127},\qquad
q_i=\operatorname{clip}(\operatorname{round}(x_i/s),-127,127),\qquad
\hat x_i=sq_i.
\]
One outlier sets $s$ for all coordinates and can waste resolution. Per-channel/group scales and techniques such as SmoothQuant address different aspects of this problem.

\section{SmoothQuant as an example of a method}
SmoothQuant does more than call tensors INT8. It applies a mathematically equivalent rescaling that moves quantization difficulty from activation outliers toward weights, then evaluates W8A8 inference empirically.

The lesson is methodological: quantization is transformation + calibration/configuration + kernel + evaluation.

\section{KV precision}
KV is unusual because it persists and is repeatedly read. Quantizing KV can reduce memory and bandwidth, potentially increasing batch capacity, but its error influences attention scores/values across many later decode steps.

Evaluate KV precision across context lengths and positions rather than only one short prompt.

\section{Error metrics}
For reference vector $y$ and low-precision $\hat y$ report, at minimum:
\[
e_\infty=\max_i|y_i-\hat y_i|,
\qquad
e_{\rm rel}=\frac{\|y-\hat y\|_2}{\max(\|y\|_2,\epsilon)}.
\]
Logit argmax agreement is useful but weak: two distributions can share argmax while calibration/ranking changes substantially.

\section{Quality gates}
\Plate{05-gate}{A precision policy needs numerical and task/output-quality evidence plus fallback.}{evod37:gate}
Define a gate using representative data and the actual model:
\begin{enumerate}
\item baseline outputs/metrics at trusted precision;
\item candidate low-precision outputs;
\item numerical parity distribution;
\item task/perplexity/accuracy/structured-output metrics as appropriate;
\item long-context and adversarial/outlier cases;
\item explicit acceptance thresholds;
\item fallback when unsupported.
\end{enumerate}

The reference's max-error/relative-error/argmax gate is only a teaching example.

\section{Calibration}
Post-training quantization parameters are learned/selected from calibration data or model statistics. Calibration data should represent expected activation distributions, context lengths and domains.

A tiny convenient calibration sample can pass while production outliers fail.

\section{Saturation and underflow}
Low precision can saturate large values or erase small ones. Tests should include extreme values and compare where the first material divergence enters the graph.

Silently clipping is a numerical policy and must be measured.

\section{Precision and batching interact}
Lower KV bytes/token can admit more pages/requests; larger batches can improve utilization. But larger batches may increase queueing or per-iteration duration. Precision can therefore shift the throughput--latency operating frontier rather than merely scale memory.

Measure the combined configuration.

\section{Chunked prefill and tails}
Long prefills can stall decodes if executed as monoliths. Chunking bounds how much prompt work joins one iteration, but smaller chunks increase scheduling overhead and can reduce prefill efficiency.

Sarathi-Serve studies this latency-throughput tension through stall-free/chunked batching. Its measured results are system-specific; the general lesson is to report both throughput and token-latency tails.

\section{Reference scheduler}
\Listing{reference}
The executable scheduler uses abstract token-work units. It prioritizes one decode token for each ready decode request, then spends remaining budget on bounded prompt chunks. The simulation records first-token and subsequent-token iteration indices.

These ticks are not GPU milliseconds.

\section{Reference precision oracle}
The same file implements FP16 conversion, a BF16-like float32 rounding transform and symmetric INT8 quantize/dequantize. The test suite checks bounded error, zero behavior and a negative case where small logit perturbations flip argmax.

That negative control matters: low RMSE does not guarantee the same discrete output.

\section{SLO experiment design}
A production evaluation should define:
\[
\mathcal W=(\lambda,\mathcal D_P,\mathcal D_O,\text{arrival process},\text{tenant mix}),
\]
where $\lambda$ is load and $\mathcal D_P,\mathcal D_O$ prompt/output-length distributions.

For each configuration report throughput, TTFT p50/p95/p99, ITL p50/p95/p99, completion latency, memory high-water mark and rejection/preemption rates.

\section{No context-free winner}
A configuration can dominate at low load and lose near saturation. Compare curves across offered load, not one cherry-picked point.

Likewise, a quantization policy can pass one benchmark and fail another domain or long-context regime.

\section{What Chapter 37 does not establish}
It does not supply a production fairness scheduler, native FP8/INT8 kernels or a universal quality threshold. It establishes the experiment/control plane required before such claims.

Speculative decoding and attention-kernel fusion remain Chapter 38.

\section{Twelve worked exercises}
\begin{enumerate}
\item Why iteration-level batching? \textbf{Requests enter/leave between autoregressive steps instead of waiting for whole-batch completion.}
\item Why chunk prefill? \textbf{Bound prompt work per iteration and reduce decode stalls.}
\item Decode-first always optimal? \textbf{No; it is one SLO policy.}
\item TTFT versus ITL? \textbf{Arrival-to-first-output versus gaps between later outputs.}
\item Throughput numerator? \textbf{Must state generated tokens, total tokens or requests.}
\item BF16 advantage over FP16? \textbf{Wider exponent range; FP16 has more significand precision in-range.}
\item Does INT8 halve FP16 bytes? \textbf{Payload often does, but scales/metadata/alignment/kernel workspace remain.}
\item Why are outliers problematic for per-tensor INT8? \textbf{They set a coarse scale for ordinary values.}
\item Same argmax means same model quality? \textbf{No.}
\item Why test KV quantization at long context? \textbf{Its stored error is reused throughout later attention.}
\item Why evaluate load curves? \textbf{Scheduling behavior changes near saturation.}
\item Release rubric? \textbf{Workload declaration; token/page budgets; TTFT/ITL tails; fairness policy; typed precision; calibration provenance; numerical+task gates; fallback; load curves; baseline comparison.}
\end{enumerate}

\section{Bridge}
Chapter 38 adds two final engine optimizations: speculative draft/verify execution and optimized attention kernels. Both must preserve the causal/paged/precision contracts developed in Chapters 35--37 and earn speed claims through synchronized hardware measurement.
